\subsection{II : algèbres graduées et degrés}
\label{note:base2-graded}
La propriété universelle de l'algèbre symétrique en 1.7.4 a pour
but une algèbre commutative $B$. Dans 2.3.2, $f^rx$ appartient
à l'idéal homogène $J$ pour tout $r\geq0$ ; c'est l'inégalité
$m+rd\geq n$ qui demande $r$ assez grand et autorise alors
le passage à l'idéal premier. La somme de 3.1.8 comporte
$\mathcal O_Y$ en degré zéro et les $S_n$ pour $n>0$.
Dans 3.2.6, $\lambda,\mu$ sont des homomorphismes ; la construction
ne donne pas à elle seule des isomorphismes. Le module de 3.3.2
est $M$.

La relation tensorielle de 4.2.1 est
$(x/1)\otimes1=(f/1)\otimes\omega(x/f)$.
Les puissances amples de 4.4.1 sont des puissances d'exposant
strictement positif. Les radicaux de 8.2.10 se prennent dans
les anneaux gradués unitaires $S$ et $S^\natural$,
et non seulement dans leur partie de degré positif.
Le terme de l'algèbre $T$ en 8.2.14 est $T_n$ ; la somme homogène
de 8.7.4 utilise $J_n$ dans son terme d'indice $n$.
Le morphisme de séries formelles en 8.10.1 est déterminé par
$T\mapsto t$.
En 8.13.3, le faisceau associé au module $N$ est $\widetilde N$.
Dans 8.14.7, la composante du morphisme est
$\alpha_{nd}(f^n)$ : si $f$ est de degré $d$, sa puissance
$n$-ième est de degré $nd$.

\subsection{II, 4--5 : fibrés projectifs et finitude locale}
\label{note:base2-projective}
Pour la décomposition du groupe de Picard en 4.2.7, la base est
connexe non vide et le module localement libre $\mathcal E$ a
rang au moins deux. Sur chaque fibre de dimension positive,
$\mathcal O(m)$ détermine l'entier $m$ ; la connexité le rend
constant. En rang un, $\mathbf P(\mathcal E)\simeq Y$ et
$\mathcal O(1)$ provient déjà de $Y$, si bien que l'entier ne serait
pas unique. Sur la base vide, il n'est pas non plus déterminé.

Avec $\mathcal L=p^*\mathcal L'\otimes\mathcal O(m)$,
le calcul de III, 2.1.15 et la formule de projection donnent
$p_*\mathcal L=\mathcal L'\otimes S_m(\mathcal E)$ pour $m\geq0$,
et zéro pour $m<0$. Le morphisme adjoint $\psi^\sharp$ est un
morphisme de modules sur $\mathbf P(\mathcal E)$.
Dans 4.4.6 et 4.6.8, les prolongements utilisés sont
quasi-cohérents de type fini ; la cohérence demanderait des
hypothèses supplémentaires sur l'anneau de base.
Le schéma projectif de 4.5.2 est $P=\operatorname{Proj}(S)$.
L'ouvert de 4.5.3 est quasi-compact, ce qui permet un choix fini
dans l'argument. Les ouverts de 5.6.1 portent l'indice $U''_k$
fixé par leur définition.

\subsection{II, 5.6 et 8.11 : propreté et contraction}
\label{note:base2-proper}
Dans 5.6.3, la factorisation de $j$ est
$j=q_2\circ(f\times1)\circ\Gamma_j$.
Le graphe est une immersion fermée parce que $P$ est séparé sur
$S$ ; cette factorisation permet d'utiliser les propriétés de
stabilité des morphismes propres.

La proposition 8.11.7 porte sur une contraction
$q:\mathbf V(\mathcal L)\to C$ propre et qui est un isomorphisme
hors de la section nulle, vers $C-\varepsilon(Y)$.
Ces conditions sont celles qui permettent de comparer cette
contraction au spectre de l'algèbre des sections des puissances
de $\mathcal L$. La seule commutativité du diagramme ne les
remplace pas : elle autoriserait par exemple une application
entièrement contenue dans la section $\varepsilon(Y)$.

\subsection{II, 8.3.9 : somme directe de faisceaux}
\label{note:base2-sheafsum}
Une section d'une somme directe de faisceaux a localement un
nombre fini de composantes non nulles. Cette finitude ne doit
pas être imposée globalement sur un ouvert arbitraire.
Sur la réunion disjointe d'une infinité de points, on peut
prendre une section dont la $n$-ième composante vaut un au
$n$-ième point et zéro ailleurs. Au voisinage de chaque point
une seule composante intervient, mais aucune famille finie
d'indices ne suffit sur l'espace entier.
La condition de 8.3.9 se formule donc après restriction à un
voisinage de chaque point de $U$.

\subsection{II : domaines des morphismes}
\label{note:base2-domains}
En 7.2.1, les points correspondants vérifient $f(x')=y'$.
Le morphisme de cônes en 8.7.4 va de $C_X$ vers $C''_X$ ;
les deux lignes du diagramme ont cette même direction.
Dans l'addendum à I, 1.8.3, le morphisme est
$Y\to\operatorname{Spec}(A)$. Dans l'index des notations,
$f$ est une section de la $A$-algèbre $B$, ce qui fixe
l'anneau dans lequel elle est évaluée.

\subsection{II : renvois}
\label{note:base2-refs}
Les résultats invoqués en 2.3.2, 3.1.7 et 4.4.6 sont
respectivement 2.1.8, 3.1.4 et I, 9.5.10.
La formule de 8.3.9 est 8.2.6.2 ; le renvoi de l'addendum
I, 1.8.10 est 1.8.9.1. L'entrée de l'index sur les algèbres
graduées se rapporte à 8.3.1.
\par
